\subsection{Rancangan Client}
\label{subsec:rancangan-client}

Rancangan client tidak menjadi bagian dari rancangan solusi utama, melainkan menjadi implikasi dari rancangan solusi yang telah diuraikan pada subbab sebelumnya.

Pada DecMed, terdapat tiga \textit{client} utama, yaitu \textit{Patient Client}, \textit{Hospital Client}, dan \textit{Ministry Client}. Perubahan pada rancangan \textit{client} hanya dibutuhkan pada \textit{Patient Client} dan \textit{Hospital Client}.

\subsubsection{Rancangan \textit{Patient Client}}
\label{subsubsec:rancangan-patient-client}

Perubahan pada \textit{Patient Client} dibutuhkan pada halaman berikut:

\begin{itemize}
	\item Halaman Home (List RME) disesuaikan agar tetap menampilkan daftar RME milik pasien, namun yang dibaca adalah segmen-segmen data yang dikelompokkan.
	\item Halaman Detail RME tetap menampilkan detail RME milik pasien, namun yang dibaca adalah segmen-segmen data RME. Data dikelompokkan berdasarkan \textit{dataset category}.
	\item Halaman Audit Delegasi merupakan halaman baru yang menampilkan informasi mengenai seluruh rantai delegasi akses yang terbentuk dari token Macaroon yang telah diterbitkan oleh pasien.
	      % \item Halaman Log Akses disesuaikan agar tetap menampilkan log akses milik pasien, namun yang dibaca adalah log akses yang terkait delegasi Macaroons.
	\item Halaman \textit{Scan} QR disesuaikan dengan menambahkan form untuk menentukan batasan akses yang akan diberikan.
\end{itemize}

\subsubsection{Rancangan \textit{Hospital Client}}
\label{subsubsec:rancangan-hospital-client}

Perubahan pada \textit{Hospital Client} dibutuhkan pada halaman berikut:

\begin{itemize}
	\item Halaman \textit{Dashboard} disesuaikan agar tetap menampilkan daftar akses yang dimiliki \textit{personnel} sesuai token yang dimiliki.
	\item Halaman Delegasi Akses merupakan halaman baru yang menampilkan informasi mengenai seluruh rantai delegasi akses yang terbentuk dari token Macaroon yang telah diterbitkan oleh pasien.
	\item Halaman Daftar RME pasien disesuaikan agar tetap menampilkan daftar RME milik pasien, namun sudah dikelompokan per episode RME.
	\item Halaman Detail RME disesuaikan agar tetap menampilkan detail RME milik pasien, namun yang dibaca adalah segmen-segmen data yang dikelompokkan.
	\item Halaman penulisan segmenRME disesuaikan agar tetap menampilkan form untuk menulis segmen RME milik pasien, namun yang ditulis adalah segmen-segmen data yang dikelompokkan.
\end{itemize}
